\ncdlesson{12}{Bioetik og fremtidens medicin: Hvor langt bør vi gå?}{生命伦理与未来医学：我们应当走多远？}
{DU3 Modul 5 (PD3) $\rightarrow$ Modul 6 (Studieprøven, B2 $\rightarrow$ C1)｜对标 NCE 4 \& 科技伦理与生命医学深度思辨}
{\textbf{伦理两难推演与警示句型}（\textit{rejse fundamentale etiske dilemmaer} · \textit{skabe præcedens for} · \textit{uoverskuelige konsekvenser} · \textit{en glidebane} · \textit{hvor sætter vi grænsen?}）· \textit{Jo..., desto...} 递进警示句式与医学伦理核心词汇}
{欢迎来到 Unit 3《伦理、制度与现代性思辨》的收官课！当 CRISPR 基因编辑能够抹去致命遗传病、AI 诊断系统能够在症状显现前数年预测疾病风险时，人类正站在“治愈疾病”与“人为设计生命”的临界点上。在 PD3 与 Studieprøven 面对生物科技或人工智能伦理题时，绝不能停留在笼统的“有好有坏”，而应熟练调用\textbf{\textit{rejse fundamentale etiske dilemmaer}}、\textbf{\textit{skabe præcedens for}}、\textbf{\textit{en glidebane}} 与 \textbf{\textit{uoverskuelige konsekvenser}} 进行层层推演，最终发出极具警醒力的追问：\textbf{\textit{Hvor sætter vi grænsen?}}}

% ---------------------------------------------------------------------
\begin{foerduleaeser}
带着下面两个问题阅读 Tekst A 和 Tekst B：
\begin{enumerate}
  \item \textit{Hvilke etiske dilemmaer opstår ifølge Tekst A, når genredigering bevæger sig fra sygdomsbehandling til forbedring af menneskelige egenskaber, og hvordan kan AI-diagnostik true patientens selvbestemmelse?}（根据 Tekst A，当基因编辑从治疗疾病滑向“改良人类性状”时会引发哪些伦理困境？AI 诊断又如何威胁患者的自主决定权？）
  \item \textit{Hvorfor advarer overlæge Sofie Nielsen i Tekst B mod at lade frygten for en glidebane føre til et totalt forbud mod ny medicinsk teknologi?}（在 Tekst B 中，主任医师索菲·尼尔森为什么警告不要让对“滑坡效应”的恐惧导致对新医疗技术的全面禁止？）
\end{enumerate}
\end{foerduleaeser}

% ---------------------------------------------------------------------
\ncdsubsec{二、Tekster（生命伦理深度专栏与临床学者回应双语篇）}

\subsection*{Tekst A：\textit{Fra helbredelse til optimering --- bioetikkens grænser}（生命伦理两难推演长文 · 约 365 词）}

\begin{tipbox}{语篇说明}
本篇以 CRISPR 基因编辑与 AI 预测诊断为切入口，展开严密的伦理风险推演。阅读时请重点留意\textbf{\color{NCDNavy}深蓝色标注的医疗伦理核心概念与法理框架}，以及\textbf{\color{NCDRed}深红色标注的伦理两难、先例滑坡与终极警示句型}。
\end{tipbox}

Fremtidens medicin rummer løfter, der for få årtier siden lignede ren science fiction. Med præcis genredigering som CRISPR kan forskere i dag gribe direkte ind i menneskets DNA for at fjerne alvorlige arvelige sygdomme, mens kunstig intelligens i diagnostikken kan afsløre skjulte mønstre i scanninger og blodprøver længe før de første symptomer. Men jo hurtigere teknologien udvikler sig, desto skarpere \dabold{rejser den fundamentale etiske dilemmaer} for hele samfundet.

Så længe nye bioteknologiske metoder udelukkende anvendes \da{i behandlingens tjeneste} til at helbrede livstruende lidelser, nyder forskningen bred moralsk opbakning. Problemet opstår i det øjeblik, skellet mellem sygdomsforebyggelse og ikke-medicinsk forbedring begynder at skride. Hvad udgør egentlig en ”defekt”, som bør rettes, og hvem skal definere det normale menneske? Hvis vi først tillader små genetiske justeringer for at ”optimere” sundhed, intelligens eller udseende, kan selv velmente indgreb \dabold{skabe præcedens for}, at menneskelige egenskaber forvandles til en handelsvare.

Bioetikere og medlemmer af Det Etiske Råd advarer derfor om, at samfundet let kan bevæge sig ud ad \dabold{en farlig glidebane}. I modsætning til traditionel medicin begrænser ændringer i kønsceller eller tidlige fostre sig nemlig ikke til én patient; de går i arv til kommende generationer og kan dermed få \dabold{uoverskuelige konsekvenser} for både den biologiske mangfoldighed og den sociale lighed. I værste fald risikerer vi et nyt genetisk klasseskel, hvor økonomisk privilegerede familier køber sig adgang til bedre biologiske forudsætninger.

Samme spænding mellem nytte og risiko viser sig i brugen af kunstig intelligens i sundhedsvæsenet. En klog og trænet algoritme kan redde liv gennem tidlig opsporing af kræft eller demens, men forudsiger samtidig statistiske sygdomsrisici, som patienten måske slet ikke ønsker at kende. Hvad sker der med \da{patientens selvbestemmelse} og retten til ikke at vide? Og hvordan forhindrer vi, at forsikringsselskaber eller arbejdsgivere får adgang til følsomme genetiske og digitale risikoprofiler, så \da{solidariteten} i det offentlige sundhedsvæsen smuldrer?

Teknologiske muligheder har en tendens til at blive til samfundsmæssige forventninger, medmindre vi griber ind i tide. Derfor må lovgivere, forskere og borgere stille det uafviselige spørgsmål: \dabold{Hvor sætter vi grænsen?} Kun gennem bindende internationale regler, gennemsigtig demokratisk debat og modet til at sige stop i tide kan vi sikre, at medicinen fortsat tjener menneskets værdighed frem for blind teknologisk optimisme.

\begin{transbox}
\small
\textbf{\color{NCDTeal}[Tekst A 中文精译]}\\
未来医学蕴含着几十年前还宛如纯粹科幻小说的承诺。借助 CRISPR 等精准基因编辑技术，研究者如今能够直接介入人类 DNA 以消除严重的遗传性疾病；与此同时，诊断领域的人工智能也能早在首批症状出现之前，就从影像扫描和血液样本中识别出隐藏的病变模式。然而，技术发展得越快，它为整个社会所引发的根本性伦理困境也就越加尖锐。

只要新的生物技术方法完全用于服务治病救人、治愈危及生命的疾苦，科学研究便能获得广泛的道德支持。问题恰恰产生于“疾病预防”与“非医疗目的的人体改良”之间的界线开始松动滑移的那一刻。究竟什么才算一种应当被修正的“缺陷”？又该由谁来定义“正常的人”？如果我们一旦允许为了“优化”健康水平、智力或外貌而进行微小基因调整，那么哪怕是出于善意的干预，也可能为“人类性状沦为商品”开创危险先例。

因此，生命伦理学家及丹麦伦理委员会成员警告称，社会极易走上一条危险的滑坡。与传统药物不同，对生殖细胞或早期胚胎所做的改动并不局限于某一位患者；它们会遗传给后代，从而对生物多样性与社会平等造成难以预料的后果。在最坏的情况下，我们甚至面临催生新型“基因阶层分化”的风险——经济优越的家庭得以花钱为子女购买更优的先天生理条件。

益处与风险之间的同样张力，也体现在医疗系统对人工智能的运用中。基于算法的诊断能够通过早期筛查癌症或痴呆症挽救生命，但与此同时，它也会预测出患者或许根本不愿知晓的概率性患病风险。届时，患者的自主决定权与“不知情权”将何去何从？我们又该如何防止保险公司或雇主获取敏感的基因与数字风险档案，从而避免公共医疗体系内的团结互助精神土崩瓦解？

除非我们及时出手干预，否则技术上的可能性往往有演变为社会普遍期待的趋势。正因如此，立法者、学者和公民必须提出那个无可回避的问题：我们究竟把边界划在哪里？唯有依靠具有约束力的国际法规、透明的民主辩论以及在关键时刻敢于叫停的勇气，我们才能确保医学继续服务于人的尊严，而非盲目的技术乐观主义。
\end{transbox}

\subsection*{Tekst B：\textit{Kronik-replik: ”Passivitet har også en etisk pris”}（临床学者回应评论 · 约 155 词）}

\textbf{阅读任务（理解与核对）：}\textit{Læs overlæge og genforsker Sofie Nielsens modsvar til debatten om bioetik, og besvar spørgsmålene nedenfor.}

\begin{emailbox}{Kronik-replik｜Fra: Sofie Nielsen (overlæge og genforsker) \quad Emne: Etik handler også om at handle i tide}
Advokater og bioetikere gør helt ret i at advare mod \dabold{en glidebane}, hvor menneskets gener bliver et kommercielt marked. Ingen seriøse forskere i Danmark ønsker designerbørn eller et samfund uden solidaritet og lige adgang.

Alligevel savner jeg ofte én afgørende vinkel i den offentlige debat: Også tøven og forbud har dybe menneskelige omkostninger. Når lovgivningen af ren frygt for \dabold{uoverskuelige konsekvenser} bremser sikre, målrettede genterapier eller lovende AI-screening, svigter vi de patienter og familier, der i dag lever med smertefulde og uhelbredelige sygdomme. At nægte at hjælpe, selvom de medicinske redskaber faktisk findes, er også et tungt moralsk valg.

Derfor bør vi ikke besvare spørgsmålet ”\dabold{hvor sætter vi grænsen?}” med et blankt nej til al medicinsk udvikling. I stedet må vi trække skarpe juridiske rammer, der klart adskiller behandling fra forbedring, kræve fuldt informeret \da{samtykke} og sikre lige adgang gennem det offentlige sundhedsvæsen.
\end{emailbox}

\textbf{Tekst B 阅读核对（Rigtigt / Forkert / Kort svar，答案见第六部分）：}
\begin{enumerate}
  \item \textit{Sofie Nielsen mener, at danske forskere bør have fri adgang til at skabe designerbørn på et kommercielt marked.} \underline{\hspace{2.2cm}}
  \item \textit{Hvilken overset etisk omkostning peger Sofie på, hvis lovgivningen bremser sikre genterapier og AI-screening?} \underline{\hspace{2.2cm}}
  \item \textit{Hvilke tre løsninger foreslår Sofie i stedet for et blankt nej til teknologisk udvikling?} \underline{\hspace{2.2cm}}
\end{enumerate}

\begin{tipbox}{本课发音与语调要点（Udtale \& Prosodi）}
\begin{itemize}
  \item \textbf{生命伦理与法律高阶借词重音}：\textit{bioetik} \ipa{biːoeˈtˢiɡ}（重音在最后音节 \textit{-tik}）；\textit{et dilemma} \ipa{diˈlɛmɑ}；\textit{en præcedens} \ipa{pʰʁɛseˈdɛnˀs}（法/拉丁借词！\textit{c} 读作 \ipa{s}，重音在最后且带 \textit{stød}）；\textit{en konsekvens} \ipa{kʰʌnsəˈkʰvɛnˀs}（重音在尾音节，带 \textit{stød}）。
  \item \textbf{复合并列与滑音细节}：\textit{en glidebane} \ipa{ˈɡliːðəˌbæːnə}（软 \textit{d}）；\textit{uoverskuelig} \ipa{ˈuːʌʊ̯ɐˌsɡuːəli}（否定前缀 \textit{u-} 常重读！\textit{over-} 读 \ipa{ʌʊ̯ɐ}，词尾 \textit{g} 不发音）；\textit{et samtykke} \ipa{ˈsɑmˌtˢyɡə}；\textit{arvemasse} \ipa{ˈɑːvəˌmɑsə}。
  \item \textbf{递进句与伦理设问语调}：朗读 \textit{Jo hurtigere..., desto skarpere...} 时，两个比较级副词（\textit{hurtigere}, \textit{skarpere}）须读出陡升强调；最后一段的 \textit{Hvor sætter vi grænsen?} 用冷静而坚定的降调收束在 \textit{\textbf{græn}sen} 上，极具警示权威感。
\end{itemize}
\end{tipbox}

\begin{transbox}
\small
\textbf{\color{NCDTeal}[Tekst B 中文对照]}\\
法学家与生命伦理学家完全有理由对人类基因沦为商业市场的“滑坡效应”提出警告。在丹麦，没有任何严肃的研究者希望出现“定制婴儿”或丧失团结互助精神与平等可及性的社会。

然而，我常常感到公共讨论中缺少一个关键视角：犹豫不决与一味禁止同样伴随着沉重的人道代价。当立法因单纯害怕难以预料的后果而叫停安全、靶向的基因疗法或前景广阔的 AI 筛查时，我们便辜负了今天正在忍受痛苦与不治之症折磨的患者和家庭。明明手中已经握有有效的医疗工具却选择拒绝施救，这本身也是一种沉重的道德抉择。

因此，我们不应以对一切医学进步一刀切说“不”，来回答“我们把边界划在哪里”这一问题。相反，我们应当划定严格的法律框架将“治病”与“增强”清晰区分开来，坚持充分的知情同意，并通过公共医疗体系保障所有人的平等可及性。
\end{transbox}

% ---------------------------------------------------------------------
\ncdsubsec{三、Gloser（核心生词表）}

\begin{tipbox}{词表重点}
注意两处极易写错的双写辅音与词形：\textit{en bioetik}（生命伦理学）和 \textit{en diagnostik}（诊断学）定指形式均须\textbf{双写 \textit{k}}（\textit{bioetikken}, \textit{diagnostikken}）；\textit{et klasseskel}（阶级鸿沟）定指与复数定指均须\textbf{双写 \textit{l}}（\textit{klasseskellet, klasseskellene}）。强变化动词 \textit{at skride}（滑动、动摇）和 \textit{at gribe ind}（出手干预）均为核心不规则变化，且 \textit{skride} 完成时使用助动词 \dabold{er}！
\end{tipbox}

\begin{glosertable}
名词 & en bioetik & \textit{(bioetikken, kun ental)} & 生命伦理学 \\
名词 & et dilemma & \textit{(dilemmaet, dilemmaer, dilemmaerne)} & 两难困境，进退维谷 \\
名词 & en præcedens & \textit{(præcedensen, kun ental)} & 先例，判例先导 \\
名词 & en glidebane & \textit{(glidebanen, glidebaner, glidebanerne)} & 滑坡（效应），危险的逐步滑落 \\
名词 & en konsekvens & \textit{(konsekvensen, konsekvenser, -erne)} & 后果，影响 \\
名词 & en arvemasse & \textit{(arvemassen, arvemasser, arvemasserne)} & 基因组，遗传物质 \\
名词 & en genredigering & \textit{(genredigeringen, -er, -erne)} & 基因编辑 \\
名词 & en diagnostik & \textit{(diagnostikken, kun ental)} & 诊断技术，诊断学 \\
名词 & en selvbestemmelse & \textit{(selvbestemmelsen, kun ental)} & 自主决定权 \\
名词 & et klasseskel & \textit{(klasseskellet, klasseskel, klasseskellene)} & 阶层鸿沟，阶级分化 \\
名词 & en solidaritet & \textit{(solidariteten, kun ental)} & 团结互助精神 \\
名词 & et samtykke & \textit{(samtykket, samtykker, samtykkerne)} & 同意，知情同意书 \\
\midrule
动词 [har] & at helbrede & \textit{(helbreder, helbredte, har helbredt)} & 治愈，治好 \\
动词 [er] & at skride & \textit{(skrider, skred, er skredet)} & 滑动，松动，偏移 \\
动词 [har] & at gribe (ind) & \textit{(griber, greb, har grebet ind)} & 干预，出手制止 \\
动词 [har] & at bremse & \textit{(bremser, bremsede, har bremset)} & 阻滞，煞停，遏制 \\
\midrule
形容词 & uoverskuelig & \textit{(uoverskueligt, uoverskuelige)} & 不可预见的，难以估量的 \\
形容词 & bindende & \textit{(bindende, 不变)} & 具有法律约束力的 \\
\end{glosertable}

% ---------------------------------------------------------------------
\ncdsubsec{四、Grammatik \& Stilistik（伦理推演警示句型、递进句式与 C1 避坑）}

\begin{grammatikbox}{语法点 1：生命伦理“两难—先例—滑坡—底线”四步推演模板}
在 PD3 与 Studieprøven 处理基因编辑、人工智能、安乐死或数据隐私等前沿伦理议题时，用下面四组核心警示搭配串联论证，极具学者风范：

\vspace{0.3em}
\renewcommand{\arraystretch}{1.24}
{\small
\begin{tabularx}{\linewidth}{@{} >{\bfseries\raggedright\arraybackslash}p{3.6cm} >{\raggedright\arraybackslash}p{4.1cm} >{\raggedright\arraybackslash}X @{}}
\toprule
\textbf{推演环节与核心搭配} & \textbf{语法接口与语义功能} & \textbf{Tekst A/B 经典示范} \\
\midrule
\textbf{1. 引发伦理困境}\newline \da{rejse fundamentale etiske dilemmaer} &
必须用动词 \textit{at rejse}（提出/引发），不可用 \textit{lave}；常配复数 \textit{etiske dilemmaer}。 &
\textit{Teknologien \da{rejser fundamentale etiske dilemmaer} for hele samfundet.} \\
\midrule
\textbf{2. 开创危险先例}\newline \da{skabe præcedens for} (+ 名词 / \textit{at}-从句) &
\textit{præcedens} 零冠词单数，固定接介词 \dabold{for}，推演“一旦开此先例将如何”。 &
\textit{Selv velmente indgreb kan \da{skabe præcedens for}, at gener bliver en vare.} \\
\midrule
\textbf{3. 滑坡与失控后果}\newline \da{en glidebane} \& \da{uoverskuelige konsekvenser} &
常用动词 \textit{bevæge sig ud ad en glidebane}；后果配 \textit{få uoverskuelige konsekvenser for...}。 &
\textit{Samfundet kan bevæge sig ud ad \da{en farlig glidebane} med \da{uoverskuelige konsekvenser}.} \\
\midrule
\textbf{4. 发出底线追问}\newline \da{Hvor sætter vi grænsen?} &
动词用 \textit{sætte}（划定边界），常以独立短问句置于结尾段首，力透纸背。 &
\textit{Lovgivere og borgere må stille spørgsmålet: \da{Hvor sætter vi grænsen?}} \\
\bottomrule
\end{tabularx}
}
\end{grammatikbox}

\begin{grammatikbox}{语法点 2：“越……越……”递进警示句式（\textit{Jo..., desto...}）与紧迫临界点}
\begin{enumerate}
  \item \textbf{\dabold{Jo + 比较级 + 主语 + 动词, desto + 比较级 + 动词 + 主语}}：\newline
    这是丹麦语高阶议论文最易丢分的语序考点！牢记“\textbf{前从后主，前不倒装、后必倒装}”：
    \begin{itemize}
      \item 前半截（\textit{Jo}-从句）：比较级放最前面，后面保持正常从句语序（\textbf{主语在前，动词在后}）：\textit{Jo hurtigere \da{teknologien udvikler} sig...}
      \item 后半截（\textit{desto}-主句）：比较级放最前面，但整个 \textit{desto + 比较级} 后\textbf{紧跟变位动词}（V2 倒装）：\textit{desto skarpere \dabold{rejser den} fundamentale dilemmaer.}
    \end{itemize}
  \item \textbf{临界点警示从句：\dabold{i det øjeblik, ...}} 与 \textbf{\dabold{medmindre ... i tide}}：\newline
    用 \textit{Problemet opstår i det øjeblik, skellet begynder at skride}（问题产生于界线开始松动的那一刻）和 \textit{medmindre vi griber ind i tide}（除非我们及时干预）制造强烈的紧迫感与行动力呼吁。
\end{enumerate}
\end{grammatikbox}

\begin{grammatikbox}{语法点 3：C1 生命伦理写作四大高频避坑点}
\begin{enumerate}
  \item \textbf{\ding{172} 切勿把 \textit{Jo..., desto...} 后半句写错语序！}：\newline
    \ding{55}\ \textit{Jo mere vi optimerer generne, desto større risikoen bliver.} $\rightarrow$ \ding{51}\ \textit{Jo mere vi optimerer generne, desto større \da{bliver risikoen}.}（注意在后半句，若比较级修饰名词 \textit{større risiko}，完整短语 \textit{desto større risiko} 在前也可能触发 \textit{opstår der} 等倒装！）
  \item \textbf{\ding{173} \textit{præcedens} 既不可数也不能加不定冠词 \textit{en}}：规范表达为 \da{skabe præcedens for...}（不写 \ding{55}\ \textit{skabe en præcedens}）。
  \item \textbf{\ding{174} 区分 \textit{at helbrede}（彻底治愈）与 \textit{at behandle}（治疗/护理）}：能根治遗传病用 \textit{helbrede}；泛指医疗过程与疗法用 \textit{behandle / behandling}。
  \item \textbf{\ding{175} \textit{at skride}（滑动/松动）完成时用 \textit{er}}：写“界线已经开始松动”时必须写 \textit{skellet \da{er begyndt} at skride} 或 \textit{grænsen \da{er skredet}}，切勿用 \textit{har skredet}。
\end{enumerate}
\end{grammatikbox}

% ---------------------------------------------------------------------
\ncdsubsec{五、Mønsterøvelser（句型与修辞操练）}

\textbf{练习 A：生命伦理核心搭配与警示句型精准填空（每空填入最准确的 1 个丹麦语词）}
\begin{enumerate}
  \item Kunstig intelligens i sundhedsvæsenet \underline{\hspace{2.2cm}} fundamentale etiske dilemmaer om patientens privatliv. (rejse，现在时)
  \item At redigere gener af kosmetiske årsager kan skabe \underline{\hspace{2.2cm}} for et nyt genetisk klasseskel. (先例)
  \item Hvis vi tillader genetisk optimering, bevæger vi os hurtigt ud ad en farlig \underline{\hspace{2.2cm}}. (滑坡)
  \item Ændringer i menneskets arvemasse kan få \underline{\hspace{2.2cm}} konsekvenser for kommende generationer. (难以预料的)
  \item Lovgivere og borgere må stille det svære spørgsmål: Hvor sætter vi \underline{\hspace{2.2cm}}? (边界，带定指词尾)
  \item Jo mere præcis genredigeringen bliver, \underline{\hspace{2.2cm}} skarpere bliver behovet for klare regler.
  \item Problemet opstår i det øjeblik, skellet mellem behandling og forbedring er \underline{\hspace{2.2cm}}. (skride，完成分词)
  \item Samfundet mister kontrollen med udviklingen, \underline{\hspace{2.2cm}} vi griber ind med bindende love i tide. (除非)
\end{enumerate}

\textbf{练习 B：C1 伦理句法纠错与推演升级（修正语序/搭配错误，或按提示升级）}
\begin{enumerate}
  \item 【\textit{Jo..., desto...} 倒装纠错】\textit{Jo hurtigere teknologien vokser, desto sværere det bliver at regulere den.} $\rightarrow$
  \item 【先例搭配与冠词纠错】\textit{Dette forslag kan lave en præcedens til kommerciel brug af fostres DNA.} $\rightarrow$
  \item 【助动词与过去分词纠错】\textit{Mange etiske eksperter frygter, at grænsen mellem helbredelse og optimering allerede har skredet.} $\rightarrow$
  \item 【核心动词搭配纠错】\textit{Ny medicinsk teknologi laver mange etiske dilemmaer og får uoverskuelig konsekvenser.} $\rightarrow$
  \item 【两难推演升级】用 \textit{i behandlingens tjeneste ... men skabe præcedens for ...} 改写：\textit{Genredigering er god til at kurere sygdom, men kan også føre til designerbørn.} $\rightarrow$
  \item 【递进警示升级】用 \textit{Jo ..., desto ...} 合并：\textit{Algoritmerne bliver klogere. Det bliver også mere presserende at beskytte patientens selvbestemmelse.} $\rightarrow$
\end{enumerate}

\textbf{练习 C：中译丹实战（精准运用伦理两难推演与警示句型）}
\begin{enumerate}
  \item 人工智能在诊断中的应用引发了关于隐私保护的根本性伦理困境。
  \item 哪怕出于善意的基因干预，也可能为将人类性状商品化开创先例。
  \item 许多伦理学家警告，社会正在走上一条带来难以预料后果的滑坡。
  \item 技术发展得越快，立法者及时出手干预的需求就越迫切。
  \item 除非我们尽早划定具有约束力的法律界线，否则医疗团结互助将瓦解。
  \item 面对未来的基因医学，我们必须追问自己究竟把边界划在哪里。
\end{enumerate}

% ---------------------------------------------------------------------
\ncdsubsec{六、PD3 / Studieprøven 实战微任务与参考答案}

\begin{emailbox}{PD3 / Studieprøven 写作微任务（Etisk debatindlæg · 200--250+ ord）}
\textbf{Situation}: Et privat medicinalfirma ønsker at lancere en ny genteknologi i Danmark, der både kan fjerne alvorlige arvelige sygdomme og mod ekstra betaling reducere risikoen for almindelige livsstilssygdomme samt styrke børns fysiske udholdenhed allerede før fødslen.\\[3pt]
\textbf{Opgave}: Skriv et argumenterende debatindlæg på \textbf{200--250 ord}, hvor du analyserer de etiske muligheder og risici ved forslaget og tager klar stilling. Inddrag mindst \textbf{fire} af lektionens nøgleudtryk (\textit{rejse fundamentale etiske dilemmaer, skabe præcedens for, en glidebane, uoverskuelige konsekvenser, hvor sætter vi grænsen?, i behandlingens tjeneste}) samt en \textit{Jo..., desto...}-sætning.
\end{emailbox}

\begin{answerbox}
\textbf{Før du læser 参考回答}
\begin{enumerate}
  \item \textit{Ifølge Tekst A opstår etiske dilemmaer, når skellet mellem helbredelse og menneskelig forbedring skrider, fordi indgreb i arvemassen kan skabe præcedens for et genetisk klasseskel med uoverskuelige konsekvenser. AI-diagnostik kan desuden true patientens selvbestemmelse, retten til ikke at vide og solidariteten, hvis forsikringsselskaber får adgang til risikoprofiler.}
  \item \textit{Hun advarer mod et totalt forbud, fordi passivitet og tøven også har høje etiske omkostninger for patienter, der lider af uhelbredelige sygdomme, når sikre behandlinger faktisk findes.}
\end{enumerate}

\textbf{Tekst B 阅读核对答案}
\begin{enumerate}
  \item \da{Forkert}（她明确表示丹麦没有任何严肃研究者希望出现定制婴儿或商业基因市场：\textit{Ingen seriøse forskere i Danmark ønsker designerbørn}）
  \item \da{At man svigter patienter, der i dag lever med smertefulde og uhelbredelige sygdomme.}（辜负了正在忍受痛苦与不治之症的当代患者）
  \item \da{Skarpe juridiske rammer mellem behandling og forbedring, fuldt informeret samtykke og lige adgang i det offentlige sundhedsvæsen.}（划定区分治疗与改良的严格法律框架、坚持知情同意并确保公共医疗体系内的平等可及）
\end{enumerate}

\textbf{练习 A 答案}
\begin{enumerate}
  \item \da{rejser}（固定搭配：\textit{rejse fundamentale etiske dilemmaer} 引发根本性伦理困境）
  \item \da{præcedens}（零冠词固定搭配：\textit{skabe præcedens for} 为……开创先例）
  \item \da{glidebane}（警示隐喻：\textit{bevæge sig ud ad en farlig glidebane} 走上危险滑坡）
  \item \da{uoverskuelige}（复数形容词，修饰复数 \textit{konsekvenser}：难以预料的后果）
  \item \da{grænsen}（核心警示设问：\textit{Hvor sætter vi grænsen?}）
  \item \da{desto}（递进比较连词结构：\textit{Jo ..., desto ...}）
  \item \da{skredet}（变化动词 \textit{at skride} 过去分词，配助动词 \textit{er}）
  \item \da{medmindre}（限制条件连词：“除非……”，后接肯定从句）
\end{enumerate}

\textbf{练习 B 答案（参考）}
\begin{enumerate}
  \item \textit{Jo hurtigere teknologien vokser, desto sværere \da{bliver det} at regulere den.}（\textit{desto}-主句部分在比较级短语后必须触发 V2 倒装，将动词 \textit{bliver} 置于主语 \textit{det} 之前）
  \item \textit{Dette forslag kan \da{skabe præcedens for} kommerciel brug af fostres DNA.}（用动词 \textit{skabe} 替代 \textit{lave}，去掉多余不定冠词 \textit{en}，并将介词改用 \textit{for}）
  \item \textit{Mange etiske eksperter frygter, at grænsen mellem helbredelse og optimering allerede \da{er skredet}.}（状态变化动词 \textit{at skride} 的完成时必须使用助动词 \textit{er}）
  \item \textit{Ny medicinsk teknologi \da{rejser} mange etiske dilemmaer og får \da{uoverskuelige} konsekvenser.}（改用固定搭配动词 \textit{rejser}，并将复数形容词 \textit{uoverskuelige} 补全词尾 \textit{-e}）
  \item \textit{Genredigering bør kun bruges \da{i behandlingens tjeneste}, da kommercielle tilbud ellers kan \da{skabe præcedens for} designerbørn.}（运用课内高级名词短语与先例搭配完成论证升级）
  \item \textit{\da{Jo} klogere algoritmerne bliver, \da{desto} mere presserende \da{bliver det} at beskytte patientens selvbestemmelse.}（用 \textit{Jo..., desto...} 严格按“前不倒装、后必倒装”规范合并）
\end{enumerate}

\textbf{练习 C 答案}
\begin{enumerate}
  \item \textit{Brugen af kunstig intelligens i diagnostik rejser fundamentale etiske dilemmaer om privatlivet.}
  \item \textit{Selv velmente genetiske indgreb kan skabe præcedens for at gøre menneskelige egenskaber til varer.}
  \item \textit{Mange etikere advarer om, at samfundet bevæger sig ud ad en farlig glidebane.}
  \item \textit{Jo hurtigere teknologien udvikler sig, desto stærkere bliver behovet for lovgivning i tide.}
  \item \textit{Medmindre vi sætter bindende juridiske grænser snart, smuldrer solidariteten i sundhedsvæsenet.}
  \item \textit{Over for fremtidens genmedicin må vi spørge os selv, hvor vi sætter grænsen.}
\end{enumerate}

\textbf{PD3 / Studieprøven 高分范文（约 230 词，不含标题）}\\
\textit{\textbf{Genmedicin skal helbrede sygdom --- ikke sælge menneskelig optimering}}

\textit{Det private medicinalfirmas forslag om at kombinere sygdomsbehandling med betalt genetisk forbedring før fødslen \dabold{rejser fundamentale etiske dilemmaer} for hele samfundet. Ganske vist er det både nødvendigt og moralsk forsvarligt at bekæmpe alvorlige arvelige lidelser, når forskningen endelig stiller sikre og præcise redskaber til rådighed. Så længe genredigering udelukkende anvendes \da{i behandlingens tjeneste} til at fjerne livstruende sygdomme, tjener den direkte patienternes værdighed og livskvalitet.}

\textit{Problemet opstår imidlertid i det øjeblik, en kommerciel aktør begynder at sælge fysisk udholdenhed og forebyggende optimering mod ekstra betaling til raske familier. Her risikerer vi nemlig at \dabold{skabe præcedens for}, at menneskets arvemasse gøres til en almindelig handelsvare. Hvis velhavende forældre kan købe sig til bedre biologiske forudsætninger for deres ufødte børn, bevæger samfundet sig hurtigt ud ad \dabold{en farlig glidebane}, der kan få \dabold{uoverskuelige konsekvenser} for den sociale lighed og \da{solidariteten} i velfærdssamfundet. \da{Jo} mere markedskræfterne får lov til at definere det ”perfekte” barn, \da{desto} skarpere tegner konturerne af et nyt genetisk klasseskel sig over generationer.}

\textit{Spørgsmålet er derfor uafviseligt: \dabold{Hvor sætter vi grænsen?} Svaret bør naturligvis ikke være et totalt forbud mod genforskning, for det ville svigte patienter med alvorlige arvelige sygdomme og uhelbredelige lidelser. I stedet bør Folketinget trække en knivskarp juridisk streg: Dokumenteret terapeutisk behandling skal godkendes og finansieres i det offentlige sundhedsvæsen med fuldt samtykke, mens kommerciel genetisk forbedring af raske fostre bør forbydes ved lov.}

\small\textbf{名师点评}：正文（不含标题）约 230 词，紧扣商业化基因增强提案展开犀利而稳健的伦理辨析。首段用 \textit{Ganske vist} 与 \textit{i behandlingens tjeneste} 肯定治病价值，确立审慎立场；第二段借 \textit{i det øjeblik...}、\textit{skabe præcedens for}、\textit{en farlig glidebane}、\textit{uoverskuelige konsekvenser} 以及严丝合缝的 \textit{Jo..., desto...} 倒装句完成风险推演；末段以 \textit{Hvor sætter vi grænsen?} 设问收束，提出“公立医疗资助治病、立法禁止商业优化”的政策区分，兼具人文温度与法理严密度！
\end{answerbox}